%% ma: L12-B16-0004
%% lop: 12
%% chuong: 5
%% bai: 16
%% dang: 12.16.3
%% loai: DS
%% muc: [TH, VD, VD, VD]
%% dap_an: "ĐSSĐ"
%% nguon: "(NBV)-ĐỀ SỐ 15-MỖI NGÀY 1 ĐỀ THI 2025.pdf (D15, MỖI NGÀY 1 ĐỀ - Đề số 15), Phần II Câu 2"
%% kien_thuc: "Bài 14 (phương trình mặt phẳng), Bài 15 (phương trình đường thẳng), Bài 16 (góc giữa đường thẳng và mặt phẳng); phép đối xứng qua mặt phẳng tọa độ"
%% trang_thai: chinh-thuc
%% ghi_chu: "Đã duyệt 16/10/2026. Đáp án gốc ĐSSĐ đúng (giải độc lập, kiểm bằng sympy/numpy). Giả thiết các điểm M thuộc OC, A thuộc CN, EM vuông góc OC được ghi tường minh vì hình gốc thể hiện bằng nét vẽ. Hình dựng lại theo số liệu đề, phép chiếu song song."
\begin{ex}
Nguyên lí phản xạ ánh sáng: khi tia sáng bị phản xạ, tia phản xạ nằm trong mặt phẳng chứa tia tới và pháp tuyến của gương tại điểm tới, góc phản xạ bằng góc tới. Chú ý: khi lấy đối xứng tia $SI$ qua gương thì tia $IS'$ là tia đối của tia $IR$.

Trong căn phòng có dạng hình hộp chữ nhật, có gương phẳng tròn tâm $E$ với độ dày không đáng kể, được treo trên tường là mặt phẳng $(FOC)$. Biết rằng $OCGF$, $OCND$, $OFKD$ là các hình chữ nhật; $OC=5\,\mathrm m$, $OD=OF=4\,\mathrm m$, $OM=3\,\mathrm m$, $EM=1{,}5\,\mathrm m$, $CA=1\,\mathrm m$ ($M\in OC$, $A\in CN$, $EM\perp OC$). Một tia sáng được chiếu từ vị trí $A$ tới $E$, tia sáng phản xạ qua gương tới điểm $B$ trên tường $(OFKD)$. Chọn hệ trục tọa độ $Oxyz$ như hình vẽ, mỗi đơn vị có độ dài $1\,\mathrm m$.
\begin{center}
\begin{tikzpicture}
  \begin{scope}[xshift=-4.0cm]
    \draw[pattern=north east lines,draw=none] (-1.5,-.25) rectangle (1.5,0);
    \draw (-1.5,0)--(1.5,0);
    \draw[dashed] (0,0)--(0,2.1) node[above,font=\small]{$N$};
    \draw[->] (-1.4,1.7)--(-.05,.05);
    \draw[->] (.05,.05)--(1.4,1.7);
    \node[above left,font=\small] at (-1.4,1.7) {$S$};
    \node[above right,font=\small] at (1.4,1.7) {$R$};
    \node[below,font=\small] at (0,-.3) {$I$};
    \node[font=\small] at (-.4,.7) {$i$};
    \node[font=\small] at (.4,.7) {$i'$};
  \end{scope}
  \begin{scope}[x={(-.5cm,-.36cm)},y={(.65cm,0cm)},z={(0cm,.65cm)}]
    \coordinate (O) at (0,0,0);\coordinate (C) at (5,0,0);\coordinate (N) at (5,4,0);\coordinate (D) at (0,4,0);
    \coordinate (F) at (0,0,4);\coordinate (G) at (5,0,4);\coordinate (K) at (0,4,4);
    \coordinate (M) at (3,0,0);\coordinate (E) at (3,0,1.5);\coordinate (A) at (5,1,0);\coordinate (B) at (0,1.5,3.75);\coordinate (H) at (0,1.5,0);
    \fill[gray!12] (O)--(C)--(N)--(D)--cycle;
    \fill[blue!6] (O)--(C)--(G)--(F)--cycle;
    \draw (O)--(C)--(N)--(D)--cycle;
    \draw (O)--(C)--(G)--(F)--cycle;
    \draw (O)--(D)--(K)--(F)--cycle;
    \draw[khuat] (E)--(M);
    \draw[khuat] (B)--(H);
    \draw[thick,->] (A)--(E);
    \draw[thick,->] (E)--(B);
    \fill[gray!50,opacity=.8] plot[domain=0:360,samples=40] ({3+.4*cos(\x)},0,{1.5+.4*sin(\x)});
    \foreach \p/\d in {O/below right,C/left,G/left,F/above left,K/above,D/below right,N/right,M/below,A/below,B/above,H/below}
      \fill (\p) circle (1pt) node[\d,font=\footnotesize] {$\p$};
    \fill (E) circle (1pt) node[above left,font=\footnotesize] {$E$};
    \draw[truc] (F)--(0,0,4.9) node[above,font=\small]{$z$};
    \draw[truc] (D)--(0,5,0) node[right,font=\small]{$y$};
    \draw[truc] (C)--(5.8,0,0) node[below,font=\small]{$x$};
  \end{scope}\end{tikzpicture}
\end{center}
\choiceTF
{\True Phương trình mặt phẳng $(ACE)$ là $3x+4z-15=0$.}
{Đường thẳng $AE$ đi qua điểm $J(1;3;3)$.}
{Gọi $(d)$ là đường thẳng đối xứng với đường thẳng $AE$ qua mặt phẳng $(OCGF)$. Góc giữa đường thẳng $(d)$ và mặt phẳng $(OCGF)$ là $41^\circ$ (kết quả làm tròn đến hàng đơn vị).}
{\True Gọi tọa độ điểm $B$ là $B(x_0;y_0;z_0)$ thì $2y_0+4z_0=18$.}
\loigiai{Từ hình vẽ: $O(0;0;0)$, $C(5;0;0)$, $D(0;4;0)$, $F(0;0;4)$ và $M(3;0;0)$, $E(3;0;1{,}5)$, $A(5;1;0)$.

Mặt phẳng $(OCGF)$ là $(xOz)$: $y=0$.

a) $\overrightarrow{CA}=(0;1;0)$, $\overrightarrow{CE}=(-2;0;1{,}5)$, $[\overrightarrow{CA},\overrightarrow{CE}]=(1{,}5;0;2)$ nên $(ACE)$: $3x+4z-15=0$. Đúng.

b) $\overrightarrow{AE}=(-2;-1;1{,}5)$. Đường thẳng $AE$: $x=5-2t$, $y=1-t$, $z=1{,}5t$. Với $x=1$ thì $t=2$, $y=-1\ne3$ nên $J\notin AE$. Sai.

c) Đối xứng qua $y=0$: $A'(5;-1;0)$ và $E$ giữ nguyên, nên $(d)$ có vectơ chỉ phương $\overrightarrow{EA'}=(2;-1;-1{,}5)$. Vectơ pháp tuyến của $(xOz)$ là $\vec j=(0;1;0)$: $\sin\varphi=\dfrac{1}{\sqrt{4+1+2{,}25}}\approx0{,}371$, $\varphi\approx22^\circ\ne41^\circ$. Sai.

d) Tia phản xạ xuất phát từ $E$ có vectơ chỉ phương $(-2;1;1{,}5)$ (đối xứng của $\overrightarrow{AE}$ qua $y=0$): $x=3-2t$, $y=t$, $z=1{,}5+1{,}5t$. Tia gặp tường $(OFKD)$: $x=0$ tại $t=1{,}5$, nên $B(0;1{,}5;3{,}75)$, $2y_0+4z_0=3+15=18$. Đúng.}
\end{ex}
